\subsection{Combined-dataset area-source proxies}
\label{sec:combined-area-proxies}

The generic proxies described above are sufficient when one spatial
driver dominates the area-source component. Several GNFR sectors,
however, combine source types that are physically distinct but still
reported together in the atmospheric inventory. Waste emissions include
solid-waste handling, wastewater treatment, and residual diffuse sources.
Industrial and fugitive emissions combine many production branches whose
spatial supports differ. Offroad transport includes railways, pipeline
transport, and non-road mobile machinery. For these sectors, a single
land-cover or population layer would obscure the process structure that
is already present in the reported inventory.

The combined-dataset method therefore first builds one spatial indicator
for each interpretable source family. These indicators are then combined
with country- and pollutant-specific allocation coefficients derived from
reported emissions. For a source family $g$, these coefficients are
computed from the reported subgroup emissions as
\begin{equation}
  \alpha_{g,p,c} =
  \frac{\displaystyle\sum_{s \in g} E_{s,p,c}^{\mathrm{rep}}}
       {\displaystyle\sum_{g'}\sum_{s \in g'} E_{s,p,c}^{\mathrm{rep}}},
  \label{eq:combined-alpha-aggregation}
\end{equation}
where $E_{s,p,c}^{\mathrm{rep}}$ is the reported emission of subgroup
$s$, pollutant $p$, and country $c$. This sector-family recombination is
introduced here as part of the upgraded proxy workflow. It builds on the
established foundation of high-resolution emission downscaling, in which
inventory totals are preserved while ancillary spatial data are used to
redistribute emissions within the coarse grid
\citep{UrbEm2021,NAEI2024Spatial}. In the present workflow, the final
step always remains the CAMS-cell normalisation introduced in
Section~\ref{sec:class-extent}; the combined proxy changes only the
within-cell pattern, not the CAMS-cell emission total.

For a sector with source families $g$, the raw score for pollutant $p$ is
written as
\begin{equation}
  S_{p,j} =
  \sum_g \alpha_{g,p,c(j)} P_{g,j},
  \label{eq:combined-sector-score}
\end{equation}
where $P_{g,j}$ is the spatial indicator for family $g$ at fine-grid
cell $j$, and $\alpha_{g,p,c(j)}$ is the reported contribution of that
family to pollutant $p$ in the country $c$ containing cell $j$, computed
as in Equation~\ref{eq:combined-alpha-aggregation}. The final weight is then
\begin{equation}
  W_{p,j \mid C} =
  \frac{S_{p,j}}{\sum_{k \in C} S_{p,k}},
  \qquad j \in C ,
  \label{eq:combined-sector-normalisation}
\end{equation}
with a uniform fallback only when all candidate indicators are zero in a
CAMS cell. This structure keeps the sectoral mass balance explicit while
allowing different pollutants to follow different spatial patterns when
their reported source composition differs.

\subsubsection{Waste}
\label{sec:combined-waste}

Waste is treated as a combined proxy because the sector brings together
solid-waste disposal, wastewater treatment, and residual sources that do
not share the same spatial footprint. This separation is important for
methane and nitrous oxide: wastewater emissions depend on collection,
treatment technology, treatment objectives, and plant size, while
solid-waste emissions are tied to disposal and handling sites
\citep{moore_naturewater2025,Song2024WastewaterNetZero}. The proxy
therefore uses three families: solid waste, wastewater, and residual
diffuse waste.

\begin{table}[htbp]
  \centering
  \caption{Spatial evidence used for the combined waste area-source proxy.}
  \label{tab:waste-spatial-rules}
  \scriptsize
  \begin{tabular}{p{0.18\textwidth} p{0.25\textwidth} p{0.45\textwidth}}
    \toprule
    Family & Spatial weighting & Retained evidence \\
    \midrule
    Solid waste &
    $1.00$ dump-site support, $0.35$ industrial or commercial support, and $0.50$ mapped solid-waste support. &
    CORINE 132 dump sites; CORINE 121 industrial or commercial units; OSM landfill, waste\_disposal, recycling, and other solid-waste-related mapped polygons or buffered points. \\
    Wastewater &
    $0.45$ agglomeration cover, $0.35$ treatment-plant support, $0.35$ population, $0.25$ imperviousness, and $0.08$ industrial context. &
    Reported wastewater agglomeration polygons; reported treatment-plant points buffered around facilities; population density; impervious surface; CORINE 131 mineral extraction, 132 dump sites, and 133 construction sites as industrial context. \\
    Residual diffuse waste &
    $0.50$ population, $0.35$ rural settlement support, and $0.15$ imperviousness. &
    Population density; rural settlement classes from the settlement model; impervious surface as a settlement support indicator. \\
    \bottomrule
  \end{tabular}
\end{table}

Each family proxy is constructed from the spatial evidence listed in
Table~\ref{tab:waste-spatial-rules} after quantile-based scaling of the
input indicators. Treatment plants are buffered and slightly smoothed so
that wastewater treatment is represented as a local treatment zone rather
than as a single fine-grid cell.

\begin{table}[htbp]
  \centering
  \caption{Waste source families used for pollutant-specific allocation coefficients.\tablefootnote{Reported source-group definitions follow the EMEP/EEA air pollutant emission inventory guidebook \citep{EMEP_EEA}.}}
  \label{tab:waste-alpha-groups}
  \scriptsize
  \begin{tabular}{p{0.18\textwidth} p{0.18\textwidth} p{0.30\textwidth} p{0.26\textwidth}}
    \toprule
    Proxy family & Reported codes & Reported source groups represented & Interpretation in the area proxy \\
    \midrule
    Solid waste &
    5A; 5B1; 5B2 &
    Solid-waste disposal on land and biological treatment of waste. &
    Emissions are allocated to dump sites, waste-management land, recycling or disposal features, and related mapped waste facilities. \\
    Wastewater &
    5D1; 5D2; 5D3 &
    Domestic wastewater handling, industrial wastewater handling, and other wastewater-related releases. &
    Emissions are allocated to wastewater agglomerations, treatment plants, and the serviced urban population context. \\
    Residual diffuse waste &
    5C2; 5E &
    Small combustion of waste, other waste handling, and waste activities not represented by the two facility-oriented families. &
    Emissions are allocated more diffusely using population, rural settlement, and impervious-surface indicators. \\
    \bottomrule
  \end{tabular}
\end{table}

Waste incineration codes 5C1a and 5C1b are excluded from the area-source
family recombination because they are not treated as diffuse residual
waste in this proxy.

For each pollutant, the three waste families are combined as
\begin{equation}
  S_{\mathrm{waste},p,j} =
  \alpha_{\mathrm{solid},p,c(j)} P_{\mathrm{solid},j}
  + \alpha_{\mathrm{ww},p,c(j)} P_{\mathrm{ww},j}
  + \alpha_{\mathrm{res},p,c(j)} P_{\mathrm{res},j}.
  \label{eq:waste-composite}
\end{equation}
If one family has no local support in a given fine cell, its coefficient
is not allowed to create emissions there. The active family coefficients
are renormalised locally before the CAMS-cell normalisation in
Equation~\ref{eq:combined-sector-normalisation}. This avoids suppressing
the proxy when, for example, wastewater facilities are absent in a rural
cell but the residual waste component remains plausible.


The detailed visual workflow for the waste proxy construction is given
in Appendix~\ref{app:ceip-weights}, 
\subsubsection{Industry}
\label{sec:combined-industry}
Industrial area emissions are represented through four source families, each targeting
a distinct production environment: refineries and petroleum handling \textit{G1}, general
manufacturing combustion and residual industry \textit{G2}, mineral products \textit{G3}, and
chemicals and metals \textit{G4}.  The four-group decomposition is introduced here because the pollutant
composition and spatial footprint of these activities differ substantially. Refinery
stacks and petroleum storage are geographically concentrated and emit a profile
dominated by $\text{SO}_x$ and VOCs, whereas general manufacturing combustion produces a
broader $\text{NO}_x$ and PM mixture distributed across generic industrial zones. Mineral
production --- cement, lime, glass, and quarrying --- generates process PM and $\text{SO}_x$
from sites that are spatially distinct from general industry, and is therefore supported
by CORINE class~131 (mineral extraction sites) combined with OSM quarry and works
features. Chemical manufacturing and primary metal production are strong point-like
sources of $\text{SO}_x$, heavy metals, and reactive nitrogen whose spatial concentration
differs markedly from both general manufacturing and mineral sites
\citep{EMEP_EEA, CAMS-REG-v4}. Treating all four families under a single generic
industrial proxy would misallocate these contrasted emission profiles across the grid.
Each group is assigned country- and pollutant-specific coefficients
$\alpha_{g,p,c}$ derived from CEIP-reported subgroup emissions following
Eq.~\ref{eq:combined-alpha-aggregation}. G1 through G4 represent the residual
area mass, which for petroleum handling likely corresponds to smaller storage and
distribution infrastructure not individually reported at high spatial precision
\citep{Crippa2016, CAMS-REG-v4}.
For each industrial group $g$, the raw spatial indicator is constructed from a
group-specific OSM component, a group-specific CORINE component, and a weak population
fallback:
%
\begin{equation}
  p_{g,j} =
  \begin{cases}
    w_{g,\mathrm{sec}}
    \bigl(
      w_{g,\mathrm{osm}}\, \mathcal{Z}(O_{g,j})
      + w_{g,\mathrm{clc}}\, \mathcal{Z}(C_{g,j})
    \bigr)
    + w_{g,\mathrm{pop}}\, P_j
    & \text{if } O_{g,j} + C_{g,j} > 0, \\[6pt]
    P_j
    & \text{otherwise,}
  \end{cases}
  \label{eq:industry-group-proxy}
\end{equation}
%
where $O_{g,j}$ and $C_{g,j}$ are respectively the OSM coverage fraction and the
CORINE score for group $g$ at pixel $j$, $P_j$ is the population proxy raster,
$\mathcal{Z}(\cdot)$ is a standardisation operator that places both inputs on a
comparable scale, and $w_{g,\mathrm{sec}} = 1 - w_{g,\mathrm{pop}}$. When no
industrial evidence is present in a pixel, the proxy falls back entirely to population,
consistent with the residual mass allocation used in \mbox{CAMS-REG-ANT} itself. The
resulting $p_{g,j}$ is subsequently normalised within each CAMS cell to obtain the
final redistribution weights.

The OSM component concentrates emissions around explicitly mapped industrial
infrastructure, while the CORINE component provides areal support where OSM tagging is
incomplete. The population term is not interpreted as a direct driver of industrial
activity; it serves only as a spatial fallback for diffuse residual emissions whose
facility evidence is locally absent. Group-specific weights are summarised in
Table~\ref{tab:industry-spatial-rules}.

For G1, the spatial indicator combines OSM features tagged as oil or fuel works,
refineries, and fuel storage tanks ($w_{\mathrm{osm}} = 0.82$) with CORINE class~121
($w_{\mathrm{clc}} = 0.15$) and a minimal population fallback
($w_{\mathrm{pop}} = 0.03$). Refineries are strongly point-like facilities; the
population term is retained only to cover residual storage and distribution
infrastructure not individually mapped. For G2, generic OSM industrial landuse,
factory and warehouse buildings, and fossil-fuel power plant features
($w_{\mathrm{osm}} = 0.55$) are combined with CORINE~121 ($w_{\mathrm{clc}} = 0.35$)
and a population fallback ($w_{\mathrm{pop}} = 0.10$); this group acts as a catch-all
for the heterogeneous residual of manufacturing combustion and ancillary industrial
processes (NFR~2D--2L), where diffuse emissions justify a larger population share. For
G3, OSM quarry, mineshaft, and mineral works features ($w_{\mathrm{osm}} = 0.65$) are
combined with CORINE~131 mineral extraction sites ($w_{\mathrm{clc}} = 0.30$) and a
small population fallback ($w_{\mathrm{pop}} = 0.05$); quarrying and mineral
processing are predominantly rural activities with low residential co-location. For G4,
OSM chemical, pharmaceutical, petrochemical, and metal works features
($w_{\mathrm{osm}} = 0.75$) are combined with CORINE~121 ($w_{\mathrm{clc}} = 0.20$)
and a minimal population fallback ($w_{\mathrm{pop}} = 0.05$); heavy industry and
primary metal production are strongly concentrated and poorly correlated with
population density.

\begin{table}[htbp]
  \centering
  \caption{Spatial evidence and weights used for the industry area-source proxy.
           $w_{\mathrm{sec}} = 1 - w_{\mathrm{pop}}$ in all cases.}
  \label{tab:industry-spatial-rules}
  \scriptsize
  \begin{tabular}{p{0.18\textwidth} p{0.24\textwidth} p{0.46\textwidth}}
    \toprule
    Group & Weights & Retained spatial evidence \\
    \midrule
    G1 --- Refineries and petroleum handling &
    $w_{\mathrm{osm}}=0.82$,\newline
    $w_{\mathrm{clc}}=0.15$,\newline
    $w_{\mathrm{pop}}=0.03$ &
    OSM: petroleum works, refineries, oil and fuel industrial sites, storage tanks,
    depots.\newline
    CORINE: 121 industrial or commercial units. \\[4pt]
    G2 --- Manufacturing combustion and residual industry &
    $w_{\mathrm{osm}}=0.55$,\newline
    $w_{\mathrm{clc}}=0.35$,\newline
    $w_{\mathrm{pop}}=0.10$ &
    OSM: industrial land, factories, warehouses, works, manufacturing, textile, food,
    sawmill, engineering, machine-shop, gas, oil, and biomass plant features.\newline
    CORINE: 121 industrial or commercial units. \\[4pt]
    G3 --- Mineral products &
    $w_{\mathrm{osm}}=0.65$,\newline
    $w_{\mathrm{clc}}=0.30$,\newline
    $w_{\mathrm{pop}}=0.05$ &
    OSM: quarries, mineshafts, cement, lime, glass, ceramic, and brick works.\newline
    CORINE: 131 mineral extraction sites. \\[4pt]
    G4 --- Chemicals and metals &
    $w_{\mathrm{osm}}=0.75$,\newline
    $w_{\mathrm{clc}}=0.20$,\newline
    $w_{\mathrm{pop}}=0.05$ &
    OSM: chemical, pharmaceutical, fertiliser, petrochemical, plastic, steel, iron,
    aluminium, copper, zinc, metal, foundry, smelting, and rolling-mill features.\newline
    CORINE: 121 industrial or commercial units. \\
    \bottomrule
  \end{tabular}
\end{table}


\begin{table}[htbp]
  \centering
  \caption{Industry source groups used for pollutant-specific allocation coefficients.}
  \label{tab:industry-alpha-groups}
  \scriptsize
  \begin{tabular}{p{0.18\textwidth} p{0.22\textwidth} p{0.28\textwidth} p{0.24\textwidth}}
    \toprule
    Proxy group & Reported codes & Reported source groups represented & Interpretation in the area proxy \\
    \midrule
    Refineries and petroleum handling &
    1A1b; 1A1c; 1B1b; 1B2a5; 1B2a6; 1B2b5; 2C1b &
    Petroleum refining, petroleum storage, fuel depots, and related refinery or petroleum handling sources. &
    Allocates petroleum-related process and handling emissions to refinery, storage, depot, and industrial land supports. \\
    Manufacturing combustion and residual industry &
    1A2a--1A2f; 1A2gvii; 1A2gviii; 2D1; 2D2; 2D3a; 2D3b; 2G; 2H1--2H3; 2I--2L &
    Manufacturing combustion, machinery, food, textiles, wood, paper, product use, and residual industrial activities. &
    Allocates residual industrial combustion and diffuse manufacturing-process emissions to general industrial supports. \\
    Mineral products &
    2A1; 2A2; 2A3; 2A4a--2A4d; 2A5a--2A5c; 2A6 &
    Cement, lime, glass, ceramics, quarrying, construction products, and other mineral processes. &
    Allocates process emissions toward mineral extraction, cement, construction, and industrial land. \\
    Chemicals and metals &
    2B1--2B7; 2B8a--2B8f; 2B9a; 2B9b; 2B10a; 2B10b; 2C1--2C6; 2C7a--2C7d &
    Chemical industry, pharmaceutical and fertiliser production, iron and steel, non-ferrous metals, and foundries. &
    Allocates emissions to chemical and metal-production supports. \\
    \bottomrule
  \end{tabular}
\end{table}


\begin{figure}[htbp]
  \centering
  \fbox{\parbox{0.88\textwidth}{\centering
  Placeholder for the industry proxy figure: panels for refineries and
  petroleum handling, manufacturing combustion and residual industry,
  mineral products, chemicals and metals, and final selected-pollutant
  weight.}}
  \caption{Suggested diagnostic figure for the combined industry proxy.
  The figure should show the difference between petroleum, manufacturing,
  mineral, and chemical or metal supports before
  pollutant-specific recombination.}
  \label{fig:industry-combined-proxy-placeholder}
\end{figure}

\subsubsection{Fugitive emissions}
\label{sec:combined-fugitive}

Fugitive emissions are highly process-dependent. Methane releases from
fuel extraction, handling, and transmission can be concentrated around
energy infrastructure, while some downstream and residual activities are
more diffuse. Satellite-based studies confirm that fossil-fuel methane
emissions are spatially heterogeneous and often dominated by production
basins, coal activities, and high-emitting infrastructure
\citep{Lu2023FuelMethane}. The fugitive proxy therefore separates coal
and solid fuels, oil transport and upstream handling, storage or refining
and distribution, and gas flaring or residual fuel losses before
pollutant-specific recombination.

\begin{table}[htbp]
  \centering
  \caption{Spatial evidence used for the fugitive-emissions area-source proxy.}
  \label{tab:fugitive-spatial-rules}
  \scriptsize
  \begin{tabular}{p{0.20\textwidth} p{0.23\textwidth} p{0.45\textwidth}}
    \toprule
    Group & Spatial weighting & Retained evidence \\
    \midrule
    Coal and solid fuels &
    $\lambda_{\mathrm{OSM}}=0.60$, $\lambda_{\mathrm{CLC}}=0.40$; final group score $0.80$ sector signal and $0.20$ population. &
    Coal, mine, quarry, and industrial-land features; CORINE 131 mineral extraction and 121 industrial or commercial units. \\
    Oil transport and upstream handling &
    $\lambda_{\mathrm{OSM}}=0.60$, $\lambda_{\mathrm{CLC}}=0.40$; final group score $0.80$ sector signal and $0.20$ population. &
    Pipeline, port, depot, and oil-industrial features; CORINE 123 port areas and 121 industrial or commercial units. \\
    Storage, refining, and distribution &
    $\lambda_{\mathrm{OSM}}=0.60$, $\lambda_{\mathrm{CLC}}=0.40$; final group score $0.80$ sector signal and $0.20$ population. &
    Refineries, storage tanks, fuel storage, depots, and fuel-service features; CORINE 121 industrial or commercial units and 122 road or rail networks and associated land. \\
    Gas flaring and residual losses &
    $\lambda_{\mathrm{OSM}}=0.60$, $\lambda_{\mathrm{CLC}}=0.40$; final group score $0.80$ sector signal and $0.20$ population. &
    Flares, chimneys, power plants, generators, substations, gas-industrial features, and gas-resource features; CORINE 121 industrial or commercial units and 131 mineral extraction. \\
    \bottomrule
  \end{tabular}
\end{table}

The spatial construction mirrors the industry workflow, but the retained
features are selected to represent fuel-related activities:
\begin{equation}
  P_{g,j}^{\mathrm{fug}} =
  \mathcal{N}
  \left(
    \lambda_{g,\mathrm{osm}} O_{g,j}^{\mathrm{fug}}
    + \lambda_{g,\mathrm{clc}} C_{g,j}^{\mathrm{fug}}
    + \lambda_{g,\mathrm{pop}} P_{\mathrm{pop},j}
  \right).
  \label{eq:fugitive-group-proxy}
\end{equation}
The population component is kept weak and is used only where a reported
group cannot be represented fully by mapped infrastructure.

\begin{table}[htbp]
  \centering
  \caption{Fugitive-emission groups used for pollutant-specific allocation coefficients.}
  \label{tab:fugitive-alpha-groups}
  \scriptsize
  \begin{tabular}{p{0.19\textwidth} p{0.18\textwidth} p{0.31\textwidth} p{0.24\textwidth}}
    \toprule
    Proxy group & Reported codes & Reported source groups represented & Interpretation in the area proxy \\
    \midrule
    Coal and solid fuels &
    1B1A; 1B1B; 1B1C &
    Coal mining, post-mining handling, and solid-fuel losses. &
    Allocates emissions toward extraction land and coal-related industrial features. \\
    Oil transport and upstream handling &
    1B2AI &
    Oil transport, upstream oil handling, port transfer, and depot activity. &
    Allocates emissions to pipeline, port, depot, and oil-industrial supports. \\
    Storage, refining, and distribution &
    1B2AIV; 1B2AV &
    Petroleum storage, refineries, distribution depots, and fuel handling. &
    Allocates emissions to storage tanks, refineries, depots, and transport-associated land. \\
    Gas flaring and residual losses &
    1B2B; 1B2C; 1B2D &
    Gas flaring, venting, residual fuel losses, and associated energy infrastructure. &
    Allocates emissions to flares, chimneys, energy facilities, gas infrastructure, and industrial land. \\
    \bottomrule
  \end{tabular}
\end{table}

The combined fugitive score is
\begin{equation}
  S_{\mathrm{fug},p,j} =
  \sum_g
  \alpha_{g,p,c(j)} P_{g,j}^{\mathrm{fug}} .
  \label{eq:fugitive-composite}
\end{equation}
This is especially relevant for CH$_4$, for which oil, gas, and coal
families can dominate the reported mass, but the same framework is kept
for all pollutants in order to preserve a consistent multiband output.

\begin{figure}[htbp]
  \centering
  \fbox{\parbox{0.88\textwidth}{\centering
  Placeholder for the fugitive-emissions proxy figure: panels for coal
  and solid fuels, oil transport and upstream handling, storage or
  refining and distribution, gas flaring or residual losses, and final
  CH$_4$ weight.}}
  \caption{Suggested diagnostic figure for the fugitive-emissions proxy.
  The CH$_4$ panel is the most important diagnostic because methane is the
  dominant pollutant for many fuel-related fugitive sources.}
  \label{fig:fugitive-combined-proxy-placeholder}
\end{figure}

\subsubsection{Offroad transport}
\label{sec:combined-offroad}

The offroad transport sector is constructed from three separate activity
families: railways, pipeline transport, and non-road mobile machinery.
These families are combined because they are reported together in the
inventory category, but their spatial evidence is fundamentally
different. Rail emissions follow railway infrastructure, pipeline
transport follows pipeline corridors and associated facilities, and
non-road machinery is linked to land uses where mobile equipment is
expected to operate. National spatial-emission methodologies commonly use
such infrastructure and land-use information to allocate non-road and
rail-related emissions when direct activity data are unavailable
\citep{NAEI2024Spatial}.

\begin{table}[htbp]
  \centering
  \caption{Spatial evidence used for the offroad transport area-source proxy.}
  \label{tab:offroad-spatial-rules}
  \scriptsize
  \begin{tabular}{p{0.20\textwidth} p{0.25\textwidth} p{0.43\textwidth}}
    \toprule
    Family & Spatial weighting & Retained evidence \\
    \midrule
    Rail transport &
    Railway-line coverage buffered by 150 m; activity share given by $\alpha_{\mathrm{rail},p,c}$. &
    Active railway lines and railway land-use polygons; inactive or unsuitable railway values such as abandoned, disused, razed, and proposed are excluded. \\
    Pipeline transport &
    Pipeline coverage buffered by 75 m; the present setup uses the line proxy directly. If a facility raster is supplied, the blend is $0.80$ line signal and $0.20$ facility signal. &
    Pipeline lines, pipeline-related areas, and associated pipeline facilities where available. \\
    Non-road mobile machinery &
    $0.50$ agricultural support and $0.35$ industrial support before pollutant-specific recombination. &
    CORINE agricultural classes 211--213 and 241--243, optional permanent crops 221--223, industrial and extraction classes 121, 131, 133, and optional port areas 123. \\
    \bottomrule
  \end{tabular}
\end{table}

The rail and pipeline components are built as buffered infrastructure
coverage scores, while the non-road component is a CORINE-only activity
support. The pollutant-specific score is
\begin{equation}
  S_{\mathrm{off},p,j} =
  \alpha_{\mathrm{rail},p,c(j)} Z_{\mathrm{rail},j}
  + \alpha_{\mathrm{pipe},p,c(j)} Z_{\mathrm{pipe},j}
  + \alpha_{\mathrm{nonroad},p,c(j)} P_{\mathrm{nonroad},j}.
  \label{eq:offroad-composite}
\end{equation}
Here, $Z_{\mathrm{rail}}$ and $Z_{\mathrm{pipe}}$ are standardised
infrastructure-coverage scores, and $P_{\mathrm{nonroad}}$ is the
land-cover machinery proxy. The coefficients are pollutant-specific
because the relative contribution of railway engines, pipeline transport,
and non-road machinery differs between pollutants.

\begin{table}[htbp]
  \centering
  \caption{Offroad transport source families used for pollutant-specific allocation coefficients.}
  \label{tab:offroad-alpha-groups}
  \scriptsize
  \begin{tabular}{p{0.20\textwidth} p{0.16\textwidth} p{0.32\textwidth} p{0.24\textwidth}}
    \toprule
    Proxy family & Reported code & Reported source groups represented & Interpretation in the area proxy \\
    \midrule
    Rail transport &
    1A3c &
    Railway transport emissions. &
    Allocates emissions to active railway corridors and railway land. \\
    Pipeline transport &
    1A3ei &
    Pipeline transport and compressor or distribution-related pipeline activity. &
    Allocates emissions to pipeline corridors and related mapped infrastructure. \\
    Non-road mobile machinery &
    1A3eii &
    Mobile machinery used outside road transport, including agricultural, industrial, construction, and service equipment. &
    Allocates emissions to land-cover classes where machinery use is plausible. \\
    \bottomrule
  \end{tabular}
\end{table}

After construction of $S_{\mathrm{off},p,j}$, the score is normalised
within each CAMS cell with Equation~\ref{eq:combined-sector-normalisation}.
The output is pollutant-specific because each band receives the reported
rail, pipeline, and non-road shares for that pollutant.

\begin{figure}[htbp]
  \centering
  \fbox{\parbox{0.88\textwidth}{\centering
  Placeholder for the offroad proxy figure: panels for rail coverage,
  pipeline coverage, non-road land-cover support, and the final selected
  pollutant weight.}}
  \caption{Suggested diagnostic figure for the offroad proxy. The figure
  should make clear that rail and pipeline families follow infrastructure
  corridors, while non-road machinery is represented through land-cover
  support.}
  \label{fig:offroad-combined-proxy-placeholder}
\end{figure}
